\subsection{Protocol definition}
\label{subsec:protocol-definition}

\subsubsection{Goal and research questions definition} 
The goal of this SLR is \textit{to explore the state of the art regarding the application of AI technologies in customer churn prediction in the retail sector, focusing on the techniques employed, the data used, and the perceived effectiveness and limitations}. 

To achieve this goal, we formulated \textcolor{blue}{8 Research Qestions (RQs). These RQs} directly address the core research objectives of our SLR by analyzing the contributions, methodologies, and findings of the selected primary studies. \textcolor{blue}{The 8 RQs are defined as follows.} %\newline


\begin{comment}
\textit{Publication Space RQs (PS-RQs)}:
\begin{itemize}[label={}, itemsep=1pt, parsep=0pt, topsep=0pt]
    \item \PSRQone
    \item \PSRQtwo
    \item \PSRQthree
\end{itemize}
\end{comment}

%\vspace{1em}  % adds vertical space roughly equivalent to one line
\begin{itemize}[label={}, itemsep=1pt, parsep=0pt, topsep=0pt]
    \item \RSRQone
    \item \RSRQtwo
    \item \RSRQthree
    \item \RSRQfour
    \item \RSRQfive
    \item \RSRQsix
    \item \RSRQseven
    \item \RSRQeight
\end{itemize}

\subsubsection{Digital libraries selection} 
To identify eligible studies, we queried four premier digital libraries commonly employed in systematic literature reviews: Scopus, Web of Science (WoS), ACM, and IEEE Xplore. Scopus (Elsevier) was chosen for its expansive, interdisciplinary coverage (encompassing over 300+ subject areas) and its stringent selection of peer-reviewed works\footnote{\url{https://www.elsevier.com/products/scopus/content\#3-selection-standards}}. The WoS Core Collection (Clarivate) comprises more than 22,000 journals, applies rigorous evaluation criteria, and has only about 10–12\% acceptance rate for the Science Citation Index Expanded\footnote{\url{https://clarivate.libguides.com/librarianresources/coverage}}. ACM and IEEE are leading professional organizations that publish key journals and conference proceedings, and organize major international scientific conferences in computer science and on artificial intelligence in particular.

\subsubsection{Search string formulation} 
To systematically construct the search string for identifying primary studies relevant to our SLR, we followed the PICOC (Population, Intervention, Comparison, Outcome, Context) framework proposed by Petersen et al. \cite{PETERSEN20151}. In our case, the \textit{Population} is represented by \textit{Customers} of the retail sector. The \textit{Intervention} focuses on \textit{Artificial Intelligence} techniques, and the desired \textit{Outcome} is the \textit{Churn Prediction}. The \textit{Comparison} viewpoint does not apply to our study as we are not interested in comparing different interventions (e.g., traditional and AI-based approaches to customer churn prediction) while the \textit{Context} viewpoint is represented by the \textit{Retail Sector}, as we are interested in surveying the literature on applications of AI to predict churn of retail customers.

For each main term within the PICOC viewpoints, a list of relevant synonyms, abbreviations, and alternative spellings was identified. The search string was then constructed by combining these terms using AND and OR operators, customized to match the syntax required by the search engine of each selected digital library, and applied to the title, abstract, and keywords fields.
Following Kitchenham et al.'s \cite{Kitchenham-Charters-guidelines} guidelines, the search string was also validated using three control papers \cite{S19-ID363-Agbemadon, S16-ID325-Matuszelanski, S11-D105-Sunarya} manually selected during a preliminary pilot study.

The resulting PICOC and the search strings customized for the different digital libraries are reported in \ref{subsec:Appendix-PICOC-and-Search-String}.


\subsubsection{Inclusion and exclusion criteria definition}  
Inclusion and exclusion criteria (IC/EC) were defined to delineate the scope of the review. 
%All studies retrieved from the selected digital libraries were exported and consolidated into a single spreadsheet and manually screened against the IC/EC. 
A paper retrieved from the selected digital libraries using the defined search string will be retained only if it satisfies all IC and does not meet any EC \cite{Kitchenham-Charters-guidelines}. Specifically, studies will be excluded if they are position papers or short papers; if they are not peer-reviewed; if the full text is not accessible; if they are secondary or tertiary studies such as reviews, mini-reviews, or surveys; if they are duplicates or earlier versions of another candidate paper; or if they primarily replicate or test the methodology of other authors rather than presenting original research. 

To qualify for inclusion, a publication must be written in English, focus specifically on customer churn prediction using artificial intelligence methods, utilize datasets from the retail sector, and base its predictions on customer and/or transactional purchase data.

The IC/EC criteria are reported in detail in \ref{subsec:Appendix-IC-EC}.


\subsubsection{Quality assessment criteria definition}
To ensure that the studies included in this SLR provide methodologically sound and reliable evidence, we defined and applied a set of six Quality Criteria (QCs) following the guidelines proposed by Kitchenham et al. \cite{SEGRESS_Kitchenham_2023}. These criteria assess the methodological rigor, clarity of reporting, and validity of the findings, supporting a transparent evaluation that strengthens the credibility of the review’s conclusions.
The list of defined QCs, along with the possible scores and their associated rationales, is presented in Table \ref{tab:QC}. The overall quality of each study was determined by summing these scores. 
% Please add the following required packages to your document preamble:
% \usepackage{graphicx}
% \usepackage[normalem]{ulem}
% \useunder{\uline}{\ul}{}
% nel preambolo:

\begin{table}[H]
\centering
\caption{Quality assessment criteria}
\label{tab:QC}
\resizebox{0.9\textwidth}{!}{
\begin{tabular}{
  |>{\columncolor{lightgray}}p{0.3\textwidth} 
  |p{0.3\textwidth}
  |p{0.3\textwidth}
  |p{0.3\textwidth}
  |
}
  \cline{2-4}
  \rowcolor{lightgray}
  \multicolumn{1}{p{0.3\textwidth}|}{\cellcolor{white}} 
    & \multicolumn{3}{c|}{\raisebox{0.4\height}{\rule{0pt}{3ex}\large{\textbf{Possible scores}}}} \\
      \hline
  \rowcolor{lightgray}
  \raisebox{0.4\height}{\rule{0pt}{3ex}\large{\textbf{Quality Criteria}}}
    & \multicolumn{1}{c|}{\raisebox{0.4\height}{\rule{0pt}{3ex}{\textbf{Fully satisfied (+1.0)}}}}
    & \multicolumn{1}{c|}{\raisebox{0.4\height}{\rule{0pt}{3ex}{\textbf{Partially satisfied (+0.5)}}}}
    & \multicolumn{1}{c|}{\raisebox{0.4\height}{\rule{0pt}{3ex}{\textbf{Not satisfied (+0.0)}}}} \\
  \hline
  \textbf{QC1: Is the paper published in a qualified venue?}
    & The paper is published in a journal ranked in Q1 or Q2 quartile or in a CORE A+ or A Conference 
    & The paper is published in a venue ranked in Q3 quartile or in a CORE B Conference
    & The paper is published in a venue ranked in Q4 quartile or in a CORE C or D Conference \\
  \hline
  \textbf{QC2: Does the work clearly describe which specific AI techniques have been used to predict Customer Churn?}
    & The adopted AI technologies are described in detail (not only the specific algorithms are specified, but also details about hyperparameters and experimental settings are provided)
    & The adopted AI techniques are clearly specified but in a general way, i.e., no additional information on how they were used is provided (such as, hyperparameters setting, experimental setting, etc.)
    & No specific AI techniques are mentioned (e.g., only the use of clustering or regression methods are mentioned) \\
  \hline
  \textbf{QC3: Does the work clearly describe which dataset/variables are used to predict Customer Churn?}
    & The dataset used for churn prediction is available for download
    & The dataset used for churn prediction is described in detail but is not available for download
    & The dataset used for churn prediction is only summarily described or not described at all \\
  \hline
  \textbf{QC4: Is class imbalance dataset addressed?}
    & Class imbalance is handled with appropriate methods and its impact is clearly discussed
    & Class imbalance is mentioned and partially addressed, but with limited detail or discussion
    & Class imbalance is not mentioned or addressed \\
  \hline
  \textbf{QC5: Have the adopted AI techniques been experimentally evaluated?}
    & The performance of the adopted AI techniques has been evaluated using a rich set of appropriate metrics
    & The performance of the adopted AI techniques has been evaluated using a few metrics
    & The performance of the adopted AI has not been evaluated at all \\
  \hline
  \textbf{QC6: Are the results of the study discussed?}
    & The results are discussed in detail, with clear summaries provided
    & The results are mentioned but not thoroughly discussed or interpreted.
    & The results are not discussed or are only mentioned without meaningful analysis. \\
  \hline
\end{tabular}%
}
\end{table}



\subsubsection{Data extraction form}\nopagebreak
To facilitate the extraction of evidence from the selected sources, we designed the Data Extraction Form (DEF) reported in Table \ref{tab:DEF}. This form specifies the key pieces of evidence to be identified in the selected papers and systematically recorded to answer each of the defined RQs. To support collaboration among the authors, the DEF was implemented as an online Google spreadsheet.
\begin{table}[H]
\centering
\caption{Data extraction form.}
\label{tab:DEF}
\resizebox{0.8\textwidth}{!}{
\begin{tabular}{|>{\centering\arraybackslash}p{0.1\textwidth}|>{\raggedright\arraybackslash}p{0.9\textwidth}|}
\hline
\rowcolor{gray!20}
 \raisebox{0.2\height}{\rule{0pt}{3ex}\textbf{RQs}} & \raisebox{0.2\height}{\rule{0pt}{3ex}\textbf{Evidence to be extracted}} \\
\hline
%\multicolumn{2}{|c|}{\textbf{Publication Space RQs}}\\  
%\hline

%PS-RQ1 & \begin{tabular}[c]{@{}l@{}}Publication year \\Publication type (article, conference paper, etc.)\end{tabular}\\
%\hline
%PS-RQ2 & \begin{tabular}[c]{@{}l@{}}Publication venue (conference name or journal name) \\Areas of interest of the venue (for journals we refer to the list of subject area \\ from Scimago; for conferences we refer to the source details from Scopus)\end{tabular}\\
%\hline
%PS-RQ3 & List of the authors’ affiliations \\
%\hline
%\multicolumn{2}{|c|}{\textbf{Research Space RQs}}\\  
%\hline
\RQlabel{RQ1} & List of sentences describing the adopted definitions of “Customer Churn” \\
\hline
\RQlabel{RQ2} & List of sentences identifying data/variables/features and the most effective features used for Churn Prediction \\
\hline
\RQlabel{RQ3} & List of sentences identifying data aggregation or correlation, or combination methods (e.g., based on time periods, order total value, products’ categories, seasonality, customer segmentation, etc.) that have been used before applying AI techniques on the resulting features \\
\hline
\RQlabel{RQ4} & List of sentences identifying the AI techniques used and found to be most effective for Churn Prediction in the Retail domain \\
\hline
\RQlabel{RQ5} & List of sentences identifying the metrics used to evaluate the performance of the adopted AI techniques. \\
\hline
\RQlabel{RQ6} & All the information related to the dataset used for Churn Prediction and publicly available (publication date, URL, \# of records, \# of customers, fields included, etc.) \\
\hline
\RQlabel{RQ7} & List of sentences on identified and discussed challenges and limitations \\
\hline
\RQlabel{RQ8} & List of sentences on emerging trends and future research directions \\
\cline{1-2}
\end{tabular}
}
\end{table}
